\documentclass[UTF8,a4paper,zihao=-4,fontset=windows]{ctexart}
\usepackage[margin=22mm,headheight=15pt]{geometry}
\usepackage{xcolor,fontspec,listings,booktabs,longtable,enumitem,fancyhdr,xurl,hyperref}
\setmonofont{Consolas}[Scale=0.88]
\setCJKmonofont{Microsoft YaHei}
\definecolor{ink}{HTML}{163B54}
\definecolor{linegray}{HTML}{CDD7DE}
\lstset{basicstyle=\small\ttfamily,breaklines=true,breakatwhitespace=false,columns=fullflexible,keepspaces=true,showstringspaces=false,frame=single,rulecolor=\color{linegray},framesep=3mm,xleftmargin=3mm,xrightmargin=3mm}
\hypersetup{hidelinks,pdftitle={ModelDeck 全量离线安装说明},pdfauthor={ModelDeck 项目交付},pdfsubject={1.4.0-dev.1-offline-20260919}}
\pagestyle{plain}
\renewcommand{\arraystretch}{1.16}
\setlength{\parindent}{0pt}\setlength{\parskip}{4pt}
\setlist[itemize]{leftmargin=1.5em,itemsep=3pt,topsep=4pt}
\ctexset{section={format=\large\bfseries\color{ink},beforeskip=13pt,afterskip=6pt}}
\linespread{1.06}
\emergencystretch=3em
\begin{document}
{\LARGE\bfseries\color{ink} ModelDeck 全量离线安装说明}\par\vspace{5mm}
交付版本：1.4.0-dev.1-offline-20260919\par
目标环境：厂家 Debian 12 ARM64 / CIX P1 / 后摩 LQ50，XH2 驱动与 SDK v1.4.0\par
制作日期：2026-09-19\par
\section*{1. 阅读须知与验收状态}
本交付物包含主项目的 UI、管理 Agent、统一 API 网关、12 个 Linux ARM64 镜像、25 个权重型模型的完整文件，以及不需要独立神经网络权重的 democratic-eval-v1 票据识别服务，共 26 个注册模型。安装过程不访问外网，不拉取镜像，不下载模型，不编译项目，也不执行联网 apt/pip 安装。\par
这是完整制品的开发交付包，版本仍为 development，不是已通过批量新机验收的稳定版。制作时核验全部模型 SHA-256、镜像标签与架构、部署配置和归档完整性；目前没有空白备用设备，未完成新盒子从零安装与断网重启演练。首次正式交付前，必须按第 7 节在一台空白同型号盒子上验收并留存报告。\par
全量安装不等于全部模型同时运行。全部模型文件与镜像都会安装；默认启动 CoPaw 模式、Guard、CV、ASR、CPU/CIX TTS、票据识别和管理平台。后摩设备上的八种模型模式互斥，通过 UI 切换，不能同时强行启动。\par
本包不包含操作系统刷机镜像、内核驱动安装程序、旧盒子的密码/API Key、业务 PDF、历史日志、旧版本备份，以及另外部署的 WeKnora、数据库和 ModelRouter 项目。\par
\section*{2. 安装前必须满足的条件}
\begin{itemize}
\item 使用同型号 CIX P1 + 后摩 LQ50 盒子；本包不能用于普通 x86 服务器，也不能直接用于其他加速器。
\item 已安装厂家 Debian 12 ARM64 系统和匹配的 XH2 v1.4.0 驱动/SDK。源设备系统内核为 6.6.89-cix-build-generic，厂家系统 BUILD\_ID 为 25Q320260821。其他厂家系统版本需要先做兼容性确认。
\item \path{/dev/aipu}、\path{/dev/xh2a_ctl0} 及其他 XH2 设备节点存在；\path{/usr/local/houmo-drv-xh2_v1.4.0} 已存在。安装器不会替换驱动或固件。
\item 系统中有 \path{cix} 用户，主目录为 \path{/home/cix}，安装人员有 sudo 权限。盒子的系统登录密码由交付方提供，与本包生成的 UI 管理密码无关。
\item 厂家基础系统已具备 Python 3、curl、tar、gzip、systemd、dpkg、iptables、nftables 等基础组件。预检会列出缺失项；离线安装器不会尝试联网补装。
\item 使用空白部署环境：没有已有项目目录和 Docker 容器，不与生产服务混装。已安装 \path{docker.io} 时会拒绝安装，不会擅自卸载。包内提供 Docker CE 29.7.2、Compose 5.5.0 及对应离线 deb。
\item 安装前建议预留 250 GB 可用磁盘空间，用于原始包、解压目录、实际部署和 Docker 数据。安装器会分别检查 \path{/home/cix} 和 Docker 数据目录所在文件系统。250 GB 指可用空间，不是磁盘标称容量。
\item 管理入口使用 TCP 17780，统一 API 使用 TCP 17700。只向可信业务网络开放；不要开放 Docker socket 或内部管理代理。
\end{itemize}
如果盒子是空盘，先由厂家刷好上述系统与驱动，再执行本文。不要使用本项目安装包代替系统恢复镜像。\par
\section*{3. 交付文件及校验}
对外交付以下三个文件；安装人员可以通过移动硬盘或内网服务器先行拷贝，正式安装阶段完全离线。\par
\par\noindent\begin{minipage}{\linewidth}\begin{lstlisting}
houmo-ai-cluster-1.4.0-dev.1-offline-20260919.tar
houmo-ai-cluster-1.4.0-dev.1-offline-20260919.tar.sha256
安装说明.pdf（另附同内容的安装说明.md）
\end{lstlisting}\end{minipage}\par
最外层使用未压缩 TAR；镜像和运行环境在包内已经压缩，避免对大权重反复压缩。必须用 tar 解包，不要将其当成 ZIP。移动硬盘需要支持大于 4 GB 的单文件，例如 exFAT；FAT32 不适用。实际项目安装目录应位于盒子的 Linux 文件系统，不在 exFAT 移动盘上运行。\par
将 TAR 与 SHA256 文件放在同一目录，在盒子上执行：\par
\par\noindent\begin{minipage}{\linewidth}\begin{lstlisting}
PKG=houmo-ai-cluster-1.4.0-dev.1-offline-20260919
sha256sum -c "$PKG.tar.sha256"
\end{lstlisting}\end{minipage}\par
只有输出 \path{OK} 才能继续。校验失败应重新拷贝并复核，不得跳过校验。SHA-256 用于检出损坏，分发时还应通过可信渠道核对校验值；本包未实施数字签名。\par
Windows 交付人员也可以在 PowerShell 中执行：\par
\par\noindent\begin{minipage}{\linewidth}\begin{lstlisting}
Get-FileHash -Algorithm SHA256 .\houmo-ai-cluster-1.4.0-dev.1-offline-20260919.tar
\end{lstlisting}\end{minipage}\par
将结果与同名 \path{.sha256} 文件中的值比较。无需把制作工作目录中的 \path{build} 文件夹交给安装人员。\par
\section*{4. 解包与只读预检}
下面假设 TAR 文件位于当前目录，解压到 \path{/opt/modeldeck-delivery}。也可以选择其他有足够空间的 Linux 文件系统。\par
\par\noindent\begin{minipage}{\linewidth}\begin{lstlisting}
PKG=houmo-ai-cluster-1.4.0-dev.1-offline-20260919
sudo mkdir -p /opt/modeldeck-delivery
sudo tar -xf "$PKG.tar" -C /opt/modeldeck-delivery
cd "/opt/modeldeck-delivery/$PKG"
sudo bash scripts/install.sh --preflight-only
\end{lstlisting}\end{minipage}\par
预检会读取全包校验和，并检查架构、系统、驱动节点、cix 用户、基础依赖、端口、已有项目和磁盘空间。只有看到 \texttt{PRECHECK PASSED} 才能安装。预检不会安装软件、启动容器或修改现有项目。\par
出现已有项目/容器提示时，应停止并联系项目维护人员。本包是首次全量安装工具，不是升级器；不要删除现有项目以强行通过预检。\par
\section*{5. 执行全量离线安装}
在刚才的解压目录中执行：\par
\par\noindent\begin{minipage}{\linewidth}\begin{lstlisting}
sudo bash scripts/install.sh
\end{lstlisting}\end{minipage}\par
安装器会再次校验包，然后依次完成：离线安装 Docker、复制全部模型和项目文件、恢复固定运行环境、生成独立凭据、导入锁定镜像、注册 systemd、创建全部模型容器、启动默认服务和真实推理验收。全过程无需互联网。\par
镜像和模型体积较大，完整校验与复制需要时间。终端会显示当前阶段；不要因为暂时没有新输出就拔盘或断电。建议使用盒子的本地终端或可靠的远程终端连接。当前安装进程依赖执行它的终端，不是后台安装守护进程。\par
安装完成应看到 \texttt{INSTALLATION PASSED}。凭据不会写入安装日志，初始管理密码由每台机器随机生成，不存在统一默认管理密码。安装日志保存在 \path{/var/log/modeldeck-offline-install/}，阶段记录保存在 \path{/var/lib/modeldeck-offline-install/state.json}。\par
如果意外断电或终端中断，先恢复环境、确认没有安装进程仍在运行，再回到同一版本的解压目录执行：\par
\par\noindent\begin{minipage}{\linewidth}\begin{lstlisting}
sudo bash scripts/install.sh --resume
\end{lstlisting}\end{minipage}\par
恢复操作会重新校验包，并跳过记录为已完成的安装阶段；未完成阶段会重新执行。不要手工修改阶段记录，不要用其他版本的包恢复。若失败发生在 NPU 启动/验收阶段，应先查看运行诊断，不能反复重试或擅自复位设备。\par
\section*{6. 登录 UI 与调用统一 API}
通过 \texttt{hostname -I} 查看盒子地址。访问 \texttt{http://<盒子IP>:17780/}，使用下面命令取得初始管理密码；按页面要求登录，不要把操作系统 SSH 密码当成 UI 密码。\par
\par\noindent\begin{minipage}{\linewidth}\begin{lstlisting}
sudo cat /home/cix/gateway-admin/INITIAL_ADMIN_PASSWORD
\end{lstlisting}\end{minipage}\par
在 UI 中管理业务 API Key。初始兼容 Key 保存在 \path{/home/cix/houmo-docker-deployment/.env}，不要将该文件发送给其他人。网关地址为 \texttt{http://<盒子IP>:17700/v1}，请求使用 \texttt{Authorization: Bearer <业务Key>}。\par
下面的本机调用不把 Key 明文写入命令历史：\par
\par\noindent\begin{minipage}{\linewidth}\begin{lstlisting}
sudo bash -c '
set -a
. /home/cix/houmo-docker-deployment/.env
set +a
curl --fail --silent --show-error \
  -H "Authorization: Bearer $HOUMO_API_KEY" \
  http://127.0.0.1:17700/v1/models
'
\end{lstlisting}\end{minipage}\par
聊天类请求使用 \path{/v1/chat/completions}。票据识别使用 \path{/v1/documents/analyze?model=democratic-eval-v1}，上传原始 PDF/图片字节并设置对应 Content-Type；它不是聊天模型，不能通过聊天接口调用。具体示例见 UI 的“API 参考”。\par
全部模型已安装并不代表都显示运行中。默认 CoPaw 可用，其他后摩模式未激活是正常现象。通过 UI 的模型管理切换模式，等待健康检查与真实推理通过后再调用该模式对应模型。\par
\section*{7. 安装验收与断网重启验收}
安装器自动执行当前后摩模式的真实推理、CPU/CIX 视觉和语音模型探针，以及通过网关提交合成票据并核对 36 个选项。验收报告写入：\par
\par\noindent\begin{minipage}{\linewidth}\begin{lstlisting}
/home/cix/houmo-docker-deployment/state/offline-acceptance.json
\end{lstlisting}\end{minipage}\par
可以在交付目录重新执行默认验收：\par
\par\noindent\begin{minipage}{\linewidth}\begin{lstlisting}
sudo python3 scripts/verify_runtime.py
\end{lstlisting}\end{minipage}\par
新机首次交付前，在没有业务请求的维护窗口，执行全模式验收：\par
\par\noindent\begin{minipage}{\linewidth}\begin{lstlisting}
sudo python3 scripts/verify_runtime.py --all-modes
\end{lstlisting}\end{minipage}\par
此操作会依次切换 CoPaw、Spark、Qwen3.8、Qwen-VL、GLM-OCR、BGE-M3、Reranker 和 CosyVoice，期间各模式不可同时提供服务。全部通过后恢复测试开始时的模式。若某模式失败，立即停止后续测试，不自动复位设备、不强行恢复加载；保留报告并由维护人员处理。探针证明执行链路可用，不代表识别准确率或长期稳定性已经验收。\par
断网重启验收步骤：\par
\begin{itemize}
\item 保持盒子无互联网连接；可以保留本地管理网络。
\item 安装与验收完成后执行 \texttt{sudo reboot}，等待系统和默认服务启动。
\item 确认管理页面可访问，\texttt{sudo modelctl status} 显示配置模式与检测模式一致。
\item 回到交付目录再次运行 \texttt{sudo python3 scripts/verify\_runtime.py}，确认报告中的 \path{passed} 为 \path{true}。
\item 保存安装日志、默认/全模式验收报告和重启后报告。报告文件会被下一次验收覆盖，应在每次测试后另行复制归档。
\end{itemize}
\section*{8. 常见问题与安全处理}
\begin{longtable}{@{}p{0.29\textwidth}p{0.67\textwidth}@{}}\toprule
现象 & 处理方法 \\\midrule\endhead
校验失败、文件缺失 & 停止安装，重新拷贝大包及校验文件；不跳过检查。 \\
架构、系统或驱动不匹配 & 使用经过确认的同型号厂家系统；不要改脚本绕过检查。 \\
已有项目或 Docker 容器 & 本包不覆盖现有部署；联系维护人员制定迁移或升级方案。 \\
缺少基础系统组件 & 由系统维护人员补齐匹配的厂家基础环境；本安装器不联网补装。 \\
磁盘空间不足 & 分别检查项目目录与 Docker 数据目录所在分区；不要删除生产数据腾空间。 \\
中断后无法直接重新安装 & 使用同一个包的 \path{--resume}；不得修改阶段记录来跳过故障。 \\
镜像缺失且没有自动拉取 & 这是离线保护；复核全包校验和与镜像导入日志。 \\
NPU 隔离、模式切换失败 & 停止重复加载，保存日志，在 UI 运行诊断中按维护流程处理。 \\
票据返回 review & 需要人工复核；不等同于接口失败，不应把结果当作全部可靠。 \\
\bottomrule\end{longtable}
常用只读排查命令：\par
\par\noindent\begin{minipage}{\linewidth}\begin{lstlisting}
sudo docker ps -a
sudo modelctl status
sudo systemctl status houmo-ai ai-model-gateway gateway-admin
sudo journalctl -u houmo-ai -n 100 --no-pager
df -h /home/cix /var/lib/docker
\end{lstlisting}\end{minipage}\par
不要执行 \texttt{docker system prune}、删除模型目录、覆盖 \path{.env} 或自动重置加速器作为通用排障步骤。诊断文件对外发送前应检查是否包含凭据或业务内容。\par
\section*{9. 安装后的文件保留与交付签收}
安装并完成重启验收后，可以由维护人员归档或移走原始 TAR 和解压安装目录，以释放空间。清理前应保留本版本的安装说明、校验文件、版本清单和验收记录；之后再次运行交付验收工具需要重新解包。\par
不能清理实际运行目录 \path{/home/cix} 下的项目、模型与固定运行环境，也不能清理 Docker 数据目录。升级与自动回滚不属于本包功能；后续版本更新需要单独的升级交付流程。\par
交付签收应记录：设备型号与编号、系统/驱动版本、包的 SHA-256、安装时间、默认模式、全模式验收结果、断网重启结果和遗留问题。源盒子的历史验收结果不能替代新盒子的验收记录。\par
\end{document}
